\section{Methods}
\label{sec:sn:methods}

\subsection{Architecture template and design space}
The template has four tiers --- devices, edge nodes, gateways and cloud --- and three
application stages per request. Pre-processing runs on the edge or gateway tier,
aggregation always on the gateway tier, and analytics on the edge tier or in the cloud;
which variant is instantiated is a design decision. A compute class is a tuple
$(C,M,P^{\mathrm{idle}},P^{\mathrm{peak}},c)$ of capacity [MIPS], memory [MB], idle and
peak power [W] and normalised cost; a link class is $(B,\tau^{\mathrm{prop}},e,c)$ of
bandwidth [Mb/s], propagation delay, transfer energy [J/MB] and cost per attached node.
The gateway-to-cloud connection is fixed infrastructure ($200$\,Mb/s, $20$\,ms) and is
not a design variable. An architecture is
\begin{equation}
x=\big(n_e,\ \kappa_e,\ n_g,\ \kappa_g,\ \ell,\ \kappa_c,\ r,\ \pi\big),
\label{eq:ce:decision}
\end{equation}
with node counts $n_e,n_g$, classes $\kappa_e,\kappa_g,\kappa_c$, link class $\ell$,
replication factor $r$ of the ingest stage and placement policy $\pi$. The design space is
the Cartesian product $\X=D_1\times\dots\times D_k$, $k=8$. Four scales
(S1--S4) and three workload classes --- high-rate telemetry, deadline-dominated real-time
control and payload-dominated sensing --- are used throughout; their domains and
parameters are listed in Supplementary Tables~S8--S10.

\subsection{Evaluation model}
With $t(i)$ the tier executing stage $i$, $N_t$ its node count and $C_t$ the per-node
capacity of its class, the offered load and utilisation are
$\Theta_t=\sum_{i:t(i)=t}\Lambda w_i\mu_i$ and $\rho_t=\Theta_t/(N_tC_t)$, where
$\Lambda$ is the system arrival rate and $\mu_i=r$ for the ingest stage and $1$
otherwise. A tier is modelled as one queueing station with the aggregate capacity of its
nodes, deterministic service and Poisson arrivals, for which the Pollaczek--Khinchine
result gives the waiting term\cite{kleinrock1975queueing,harchol2013performance}, so
that
\begin{equation}
L(x)=\sum_{i=1}^{3}\frac{w_i}{C_{t(i)}}
+\sum_{i=1}^{3}\frac{w_i}{C_{t(i)}}\cdot\frac{\rho_{t(i)}}{2\,(1-\rho_{t(i)})}
+\sum_{(s,\ell)\in\mathcal{H}(x)}\Big(\frac{8s}{B_\ell}+\tau^{\mathrm{prop}}_\ell\Big),
\label{eq:ce:latency}
\end{equation}
where $\mathcal{H}(x)$ is the set of hops implied by the policy and the replication
factor. The middle term is the contention-dependent one. The aggregation is a tractable approximation of a multi-server tier rather than an exact
M/D/$c$ representation, kept deliberately because it is the coarsest form that still makes
latency depend on utilisation and keeps the bound computable in closed form. Both the discrete-event
evaluator and the physical testbed implement per-node servers, so its cost is measured
rather than assumed. Energy over a horizon $T=1$\,h combines idle and load-proportional
compute power with transfer energy; cost sums node and link coefficients; network volume
is the per-second byte flow. Feasibility requires $\rho_t\le0.85$ on every used tier,
memory and bandwidth to fit, and the latency, cost and energy budgets to hold. Full
expressions are in Supplementary Note~2.

\subsection{Admissible bounds}
Write $\C(p)$ for the completions of a partial design $p$. A bound $\LB_i$ is admissible
if $\LB_i(p)\le f_i(x)$ for all $x\in\C(p)$. Two constructions are used.

\emph{Term-wise relaxation.} If $f_i(x)=\sum_{t\in T_i}\varphi_t(x)$ with
$\varphi_t\ge0$ and $\underline{\varphi}_t(p)=\min_{x\in\C(p)}\varphi_t(x)$, then
$\LB_i(p)=\sum_t\underline{\varphi}_t(p)$ is admissible, because each term is bounded
below independently. The minima are taken per term, so the combination realising them
need not correspond to any completion; the bound is optimistic but valid. Because the decision domains are finite, the minimum of each term over the reachable
domain is precomputed exactly; for the ordered scalar domains this lookup reduces to a
domain endpoint, so the bound is evaluated in constant time.

\emph{Exactness on determined terms.} If the variables a term depends on are already
assigned, the term is constant on $\C(p)$ and its exact value may replace the relaxed
one, which preserves admissibility and tightens the bound. The decision order places the
determinants of tier utilisation --- node counts, classes, replication and policy ---
before the remaining variables, so the waiting term of \eqref{eq:ce:latency} becomes
exact early in the traversal instead of being relaxed to zero. This is what makes a
contention-dependent objective usable without assignment monotonicity. If the
determinants are fixed and $\rho_t\ge1$, no completion has finite latency and the bound
is $\infty$, which is admissible.

\subsection{Pruning rules and exactness}
Three rules are applied at each node of the prefix tree. \emph{Structural}: the
well-formedness predicate $\sigma$ (gateways do not outnumber edge nodes, replication
does not exceed the edge count, replication is undefined for policies that place ingest
off the edge tier) depends only on assigned variables and cannot be repaired by later
assignments, so a violation eliminates all completions. It defines the search space
rather than following from the constraints, and the exhaustive baseline applies the same
predicate. \emph{Feasibility}: if an admissible bound on a constrained quantity exceeds
its limit, no completion is feasible. \emph{Dominance}: if an evaluated feasible
architecture $y$ satisfies $F(y)\preceq \LB(p)$, then every completion of $p$ is
dominated by $y$ and none can be Pareto-optimal. Together with admissibility these give
front preservation: the non-dominated set returned equals the exact Pareto front of the
model in objective space. Proofs are in Supplementary Note~1. Two consequences are used
in the experiments: the guarantee is order-independent, so traversal order affects
efficiency only; and it is stated in objective space, so architectures sharing an
objective vector may be represented once.

\subsection{Verification of correctness}
Theoretical correctness and implementation correctness are checked separately. Bound
admissibility is verified exhaustively on the smallest scale --- every partial design
against every completion, for all three workloads. For every partial design that the
implementation prunes, all of its completions are enumerated and checked against the
claim of the rule that discarded it. Equality with the exhaustive front is re-checked on
S1--S3 for every pruning variant as a regression test.

\subsection{Baselines}
Exhaustive enumeration provides ground truth on S1--S3. Random search and
NSGA-II\cite{deb2002nsga2} (pymoo\cite{blank2020pymoo}, population $40$, SBX crossover,
polynomial mutation, integer repair) run at matched budgets of $100$, $250$, $500$ and
$1000$ full evaluations with $20$ seeds each; budgets are charged to actual evaluator
calls, including infeasible ones. An independent exact baseline encodes the finite
evaluated relation and solves for the non-dominated set with an SMT
solver\cite{demoura2008z3}; it requires the full table to be built first and therefore
consumes one evaluator call per structurally valid candidate. A second, compact exact
baseline avoids evaluator calls. Every decision becomes a one-hot group of Booleans, and
every additive term of the model is tabulated as exact rationals over the decisions it
depends on only --- edge utilisation, for example, over node count, node class,
replication and placement --- so that constraints and objectives are linear in the
Booleans. The non-dominated set is enumerated with the guided-improvement algorithm: find
a feasible design, add the constraint that a solution must dominate it until the system
becomes unsatisfiable, record the last vector and exclude everything it weakly dominates.
A direct polynomial encoding with utilisations and waiting terms as real unknowns was
tried first and did not complete the smallest space within five minutes, which is why the
factorised form is used. Designs returned by either solver are re-evaluated with the
ordinary evaluator only to compare binary64 objective vectors with the exhaustive front.
Recall and precision are
computed on sets of objective vectors, hypervolume\cite{zitzler1999hv} is normalised by
the exhaustive front, and IGD$^{+}$\cite{ishibuchi2015igdplus} uses that front as
reference, following the guidance on indicator choice in ref.~\citenum{audet2021indicators}.

\subsection{Discrete-event evaluator}
A higher-fidelity evaluator simulates each design with multi-server tiers, explicit link
contention, deterministic service and Poisson arrivals; it plays the role that a
discrete-event simulator such as OMNeT++\cite{varga2008omnet} plays in trace-driven
industrial methodologies\cite{saadatmand2025compdse}. It serves two purposes: as an
internal reference for the accuracy of the analytical model, and as the anchor of the
evaluation-cost analysis, at a measured $\DesCostMs$\,ms per design point averaged over
$40$ designs.

\subsection{Containerised deployment}
Each edge node, gateway and cloud instance is a separate container limited to one CPU, so
that a node behaves as the single server the model assumes at node level. Services are
written in Go and share one binary: a request carries a plan of the remaining stages,
each node consumes CPU work equal to $w_i/C_t$ and forwards the remainder with the real
payload to the next hop. Links are shaped per node with \texttt{tc netem}
\cite{hemminger2005netem} using the delay and rate of the link class under test, with a
wide-area profile on the cloud hop; shaping is applied in the forward direction of each
hop, matching the one-way transfers counted by the model. Replication is emulated
faithfully: $r$ replicas execute the ingest stage in parallel and carry the
synchronisation payload, and the remainder of the chain executes once. A load generator
produces Poisson arrivals, $900$ requests per run of which $150$ are discarded as warm-up.
Thirty-six architectures --- Pareto-optimal, dominated and near-boundary --- were selected
subject to fitting on one machine ($\le4$ edge nodes, $\le2$ gateways), and each was run
three times. The analytical model was frozen before this round.

\subsection{Statistical analyses}
Stochastic baselines are run with $20$ independent seeds per scale, workload and budget,
and are summarised by medians. Differences between NSGA-II and random search are tested
with the one-sided Mann--Whitney $U$ test\cite{mann1947whitney} on the per-seed
hypervolume ratios, with Cliff's $\delta$\cite{cliff1993delta} as effect size;
$p$-values are reported unadjusted and treated as exploratory. HCA-DSE, exhaustive
enumeration and both solver baselines are deterministic, so their counts are reported
from one run; runtimes in the cost and traversal-order experiments are medians of three
repetitions, and other runtimes are single measurements. In the deployment, each architecture
is run three times and summarised by the median of its run means. Accuracy is reported as
mean absolute percentage error and median absolute error, rank agreement as Spearman's
$\rho$, and decision agreement as the share of within-workload architecture pairs whose
measured order matches the predicted order, among pairs the model separates by more than
a given relative margin. Intervals on these shares are $95\%$ Wilson score intervals;
pairs share architectures and are not independent, so the intervals are indicative. No
observation is excluded; configurations whose repeated means differ by a factor of $1.5$
or more are flagged and retained.

\subsection{Reproducibility}
The implementation is Python 3 with NumPy, SciPy and pymoo; the deployment services are
Go under Docker Compose. Every figure and every numeric table in this paper, main and
supplementary, is generated from the raw result files by a single script, so no number is
transcribed by hand.
